\section{策略表达与生成模型}

模仿学习的关键一环是如何表达策略分布。单一的高斯回归只能表示局部、单峰行为，无法匹配多模态专家数据。

\subsection{分布表示的能力}

离散动作可以用 categorical 直接建模，而连续动作层面常见写法也包括 normalizing flow、Mixture Density Network、diffusion model。

\begin{importantbox}{生成模型带来的表达力飞跃}
将策略视为生成网络 $p_\theta(a \mid s)$，可以借鉴 diffusion, flow, VAE 等架构，让 policy 自身成为可采样、可推理的概率模型，从而自然适配多模态演示数据。
\end{importantbox}

\subsection{Diffusion、Mixture 与 Autoregressive 策略}

主流表达手段：
\begin{enumerate}
    \item \textbf{Mixture of Gaussians}：同时输出多个 $(\mu_i, \sigma_i, \alpha_i)$，用 softmax 权重混合。
    \item \textbf{Autoregressive discretization}：将动作空间拆成子维度，递归建模每一维的分布。
    \item \textbf{Diffusion Policy}：先采样噪声，再迭代去噪，最终生成动作。
\end{enumerate}

\begin{table}[H]
\centering
\begin{tabular}{p{0.22\textwidth} p{0.34\textwidth} p{0.34\textwidth}}
\toprule
方法 & 优势 & 典型场景 \\
\midrule
MoG & 提供有限且 interpretable 的混合模式 & 低维连续控制 \\
Autoregressive & 维度分解，捕捉细粒度 correlation & 自动驾驶行为序列 \\
Diffusion & 可做 implicit sampling & 高维机器人操作 \\
\bottomrule
\end{tabular}
\caption{表达性策略的典型架构}
\end{table}

\subsection{采样与覆盖诊断}

策略参数训练后需要验证它对 expert distribution 的覆盖度。常见做法包括：
\begin{itemize}
    \item KL divergence sweep：用专家数据和采样数据估计 KL。
    \item Reject sampling：在 rollout 时如果动作概率低于 threshold，触发 human override。
    \item Latent space interpolation：在 latent 控制变量中插值，观察生成动作是否落在专家 manifold。
\end{itemize}

\begin{knowledgebox}{切换采样策略}
Archit 提醒：\textit{``采样策略决定了我们是否能看到极端行为''}。在部署前应把 policy 的 temperature、top-k/top-p 调整写入 runbook 并记录 observational log。
\end{knowledgebox}

Flow-based policy 用于把动作空间映射到更规则的 latent space，训练时直接拟合 latent 的目标分布，使得采样时只需沿 latent 方向推进再通过 invertible map 还原动作。对于高维动作（如 humanoid 控制），flow 可以显著降低采样难度。

\begin{importantbox}{潜在控制的直觉}
Flow policy 等价于先在潜在空间做简单的高斯采样，再用 invertible mapping 将其转换为复杂动作。这样可以把模仿学习的目标移到 latent space 中，简化梯度传播。
\end{importantbox}

\subsection{本章小结}

策略表达决定了模仿学习是否能够覆盖 expert space。用生成模型建模多模态分布、Flow-based latent control 以及 rejection sampling，可以在 rollout 端持续纠偏，实现与 expert distribution 的近似，并为 downstream RL 微调提供稳定起点。

